\documentclass[12pt]{article}
\usepackage[margin=1in]{geometry}
\usepackage{amsmath}
\usepackage{amssymb}
\usepackage{amsthm}
\usepackage{enumitem}
\usepackage{xcolor}
\usepackage{pagecolor}
\usepackage{marginnote}
\usepackage{fancyhdr}
\usepackage{bm}
\usepackage{etoolbox}
\usepackage{mdframed}

\definecolor{cream}{HTML}{faf7f1}
\pagecolor{cream}

\newcommand{\defn}[1]{\textbf{#1}}
\newcommand{\defeq}{\mathrel{\vcenter{\baselineskip0.5ex \lineskiplimit0pt
                     \hbox{\scriptsize.}\hbox{\scriptsize.}}}=}
\newcommand{\T}{\top}
\newcommand{\F}{\bot}

\newtheoremstyle{proofstyle}
  {}{}{}{}{\bfseries}{.}{ }{}
\theoremstyle{proofstyle}
\newtheorem*{proof*}{Proof}


% Define sidenote command (simplified - uses marginnote)
\newcommand{\sidenote}[1]{\marginnote{\small\itshape #1}}

% Title formatting
\makeatletter
\renewcommand{\maketitle}{%
  \begin{flushleft}
    {\Large\bfseries CS173: Discrete Structures}\\[0.3em]
    {\Large\bfseries Problem Set 3}\\[0.5em]
    {\normalsize Mohammed Al Salhi}
  \end{flushleft}
  \vspace{1em}
}
\makeatother

\AtBeginDocument{\mathversion{bold}}

\begin{document}

\maketitle

\noindent
Throughout this problem set, we will use the following shorthand notation for axioms:
\begin{itemize}
    \item \textbf{Extensionality} refers to The Axiom of Extensionality (Axiom 1)
    \item \textbf{Pairing} refers to The Axiom of Pairing (Axiom 2)
\end{itemize}

\begin{enumerate}[left=0pt, itemsep=2em, parsep=1em]

\item Prove each of the following statements.

\begin{enumerate}[label=(\alph*), itemsep=1.5em, leftmargin=2em]

\item Prove $\forall x\forall y\forall z((x \subseteq y \wedge y \subseteq z) \Rightarrow x \subseteq z)$.

\begin{proof*}
    Let $x$, $y$, and $z$ be arbitrary sets. Assume $x \subseteq y$ and $y \subseteq z$.
    
    To prove $x \subseteq z$, we must show $\forall w(w \in x \Rightarrow w \in z)$ by the definition of subset.
    
    \begin{quote}
        Let $w$ be an arbitrary set and assume $w \in x$.
        
        Since $x \subseteq y$, we have $\forall u(u \in x \Rightarrow u \in y)$ by the definition of subset.
        
        Thus $w \in y$.
        
        Since $y \subseteq z$, we have $\forall u(u \in y \Rightarrow u \in z)$ by the definition of subset.
        
        Thus $w \in z$.
    \end{quote}
    
    Therefore $\forall w(w \in x \Rightarrow w \in z)$, which means $x \subseteq z$ by the definition of subset. \qed
\end{proof*}

\item Prove $\forall x\forall y(x \cap y \subseteq x)$.

\begin{proof*}
    Let $x$ and $y$ be arbitrary sets.
    
    To prove $x \cap y \subseteq x$, we must show $\forall w(w \in (x \cap y) \Rightarrow w \in x)$ by the definition of subset.
    
    \begin{quote}
        Let $w$ be arbitrary and assume $w \in (x \cap y)$.
        
        Since $x \cap y$ is an intersection of $x$ with $y$, we have 
        
        $\forall u(u \in x \cap y \Leftrightarrow ((u \in x) \wedge (u \in y)))$ by the definition of intersection of two sets.
        
        Thus $w \in x \wedge w \in y$.
        
        Therefore $w \in x$.
    \end{quote}
    
    Therefore $\forall w(w \in x \cap y \Rightarrow w \in x)$, which means $x \cap y \subseteq x$ by the definition of subset. \qed
\end{proof*}
\newpage

\item Prove $\forall x\forall y(x \subseteq x \cup y)$.

\begin{proof*}
    Let $x$ and $y$ be arbitrary sets.
    
    To prove $x \subseteq x \cup y$, we must show $\forall w(w \in x \Rightarrow w \in x \cup y)$ by the definition of subset.
    
    \begin{quote}
        Let $w$ be arbitrary and assume $w \in x$.
        
        Consider the definition $x \cup y := \{u \mid u \in x \vee u \in y\}$.
        
        Since $w \in x$, we have $w \in x \vee w \in y$.
        
        Thus $w \in x \cup y$.
    \end{quote}
    
    Therefore $\forall w(w \in x \Rightarrow w \in x \cup y)$, which means $x \subseteq x \cup y$ by the definition of subset. \qed
\end{proof*}
\item Prove $\forall x\forall y(x \cap y \subseteq x \cup y)$.

\begin{proof*}
    Let $x$ and $y$ be arbitrary sets.
    
    By problem 1b, we have $x \cap y \subseteq x$.
    
    By problem 1c, we have $x \subseteq x \cup y$.
    
    By problem 1a (transitivity of $\subseteq$), since $x \cap y \subseteq x$ and $x \subseteq x \cup y$, we have $x \cap y \subseteq x \cup y$. \qed
\end{proof*}

\end{enumerate}

\item Prove each of the following statements.

\begin{enumerate}[label=(\alph*), itemsep=1.5em, leftmargin=2em]

\item Prove $\forall x\forall y\exists p(p = (x, y))$.

\begin{proof*}
    Let $x$ and $y$ be arbitrary sets.
        
    By Pairing applied to $x$ and $x$, there exists a set $\{x, x\}$.
    
    By Extensionality, $\{x, x\} = \{x\}$.
    
    By Pairing applied to $x$ and $y$, there exists a set $\{x, y\}$.
    
    By Pairing applied to $\{x\}$ and $\{x, y\}$, there exists a set $\{\{x\}, \{x, y\}\}$.

   Thus, by definition, there exists an ordered pair $(x,y)$.

    Therefore $\exists p(p = (x, y))$. \qed
\end{proof*}
\newpage
\item Prove $\forall a\forall b\forall x\forall y((a, b) = (x, y) \Leftrightarrow (a = x \wedge b = y))$.

\begin{proof*}
    Let $a$, $b$, $x$, and $y$ be arbitrary sets.
    
    We prove both directions of the biconditional.
    
    \textbf{($\Rightarrow$)} Assume $(a, b) = (x, y)$.

    \begin{quote}
        By the definition of ordered pair, $\{\{a\}, \{a, b\}\} = \{\{x\}, \{x, y\}\}$.
    
        We consider two cases.
        
        \textit{Case 1:} Suppose $a = b$.
        
        \begin{quote}
            Since $a = b$, we have $\{a, b\} = \{a\}$ by Extensionality.
            
            Thus $\{\{a\}, \{a, b\}\} = \{\{a\}, \{a\}\} = \{\{a\}\}$ by Extensionality.
            
            Therefore $\{\{x\}, \{x, y\}\} = \{\{a\}\}$.
            
            By Extensionality, both $\{x\}$ and $\{x, y\}$ must equal $\{a\}$.
            
            From $\{x\} = \{a\}$, we have $x = a$.
            
            From $\{x, y\} = \{a\}$, we have $y = a$.
            
            Thus $a = x$ and $b = a = y$, so $a = x$ and $b = y$.
        \end{quote}
        
        \textit{Case 2:} Suppose $a \neq b$.
        
        \begin{quote}
            Since $a \neq b$, we have $b \in \{a, b\}$ but $b \notin \{a\}$.
            
            Thus $\{a\} \neq \{a, b\}$ by The Extensionality.
            
            Since $\{\{a\}, \{a, b\}\} = \{\{x\}, \{x, y\}\}$, we know $\{a\} \in \{\{x\}, \{x, y\}\}$.
            
            By Extensionality, either $\{a\} = \{x\}$ or $\{a\} = \{x, y\}$.
            
            Suppose $\{a\} = \{x, y\}$.
            
            \begin{quote}
                Then both $x$ and $y$ must equal $a$ by Extensionality.
                
                Thus $x = a = y$, which means $x = y$.
                
                But then $\{x, y\} = \{x\}$ by The Extensionality.
                
                Since $\{a, b\} \in \{\{a\}, \{a, b\}\} = \{\{x\}, \{x, y\}\}$ and \\ $\{x, y\} = \{x\}$, we have $\{a, b\} = \{x\}$.
                
                Combined with $\{a\} = \{x\}$, this gives $\{a, b\} = \{a\}$.
                
                This contradicts $\{a\} \neq \{a, b\}$.
            \end{quote}
            
            Therefore $\{a\} \neq \{x, y\}$, so we must have $\{a\} = \{x\}$.
            
            From $\{a\} = \{x\}$, we have $a = x$ by Extensionality.
            
            Since $\{a, b\} \in \{\{a\}, \{a, b\}\}$ and $\{a, b\} \neq \{a\} = \{x\}$, \\ we have $\{a, b\} = \{x, y\}$.
            
            From $\{a, b\} = \{x, y\}$ and $a = x$, we have $b \in \{x, y\}$.
            
            If $b = x$, then $b = x = a$, contradicting $a \neq b$.
            
            Thus $b = y$.
        \end{quote}
        
        In both cases, $a = x$ and $b = y$.
    \end{quote}
    
    
    
    \textbf{($\Leftarrow$)} Assume $a = x$ and $b = y$.
    
    \begin{quote}
        Since $a = x$, we have $\{a\} = \{x\}$ by Extensionality. 
        
        Since $a = x$ and $b = y$, we have $\{a, b\} = \{x, y\}$ by Extensionality.
        
        Therefore $\{\{a\}, \{a, b\}\} = \{\{x\}, \{x, y\}\}$ by Extensionality.
        
        By the definition of ordered pair, $(a, b) = (x, y)$. 
    \end{quote}
    


    Therefore $(a, b) = (x, y) \Leftrightarrow (a = x \wedge b = y)$. \qed
\end{proof*}

\end{enumerate}

\newpage

\item Prove each of the following statements.

\begin{enumerate}[label=(\alph*), itemsep=1.5em, leftmargin=2em]

\item Prove that $\forall x(x \cup \{x\} \neq \emptyset)$.

\begin{proof*}
    Let $x$ be an arbitrary set.
    
    By Pairing applied to $x$ and $x$, there exists a set $\{x, x\}$.
    
    By Extensionality, $\{x, x\} = \{x\}$.
    
    Since $\{x\} \subseteq x \cup \{x\}$ by Problem 1c, and $x \in \{x\}$, we have $x \in x \cup \{x\}$ by the definition of subset.    \
    
    Suppose for contradiction that $x \cup \{x\} = \emptyset$. By Extensionality, $x \cup \{x\}$ and $\emptyset$ have the same elements. But $x \in x \cup \{x\}$ and $x \notin \emptyset$, a contradiction.
    
    Therefore, $x \cup \{x\} \neq \emptyset$. \qed
\end{proof*}

\item Prove that $\forall x\forall y(x \neq y \Rightarrow x \cup \{x\} \neq y \cup \{y\})$.

\begin{proof*}
    Let $x$ and $y$ be arbitrary sets such that $x \neq y$.
    
    To prove this we will use reductio ad absurdum. 
    \begin{quote}
        Assume $x \cup \{x\} = y \cup \{y\}$.
        
        As shown in problem 3a, we have $x \in x \cup \{x\}$, and so $x \in y \cup \{y\}$.
        
        So, by Extensionality, $x \in y$ or $x \in \{y\}$. 
        
        Since $x \in \{y\}$ implies, by Extensionality, $x = y$, which contradicts $x \neq y$, we have $x \in y$.
        
        Again, as shown in problem 3a, we have $y \in y \cup \{y\}$ and so $y \in x \cup \{x\}$.

        Assume, for the sake of this proof, $y \notin x$.
                
        Since $y \in x \cup \{x\}$ and $y \notin x$, we must have $y \in \{x\}$, so $y = x$. 
        
        We have both $x \neq y$ and $y = x$, which is a contradiction.
    \end{quote}
        
    Therefore, by reductio ad absurdum, we have $x \neq y \Rightarrow x \cup \{x\} \neq y \cup \{y\}$. \qed
        
    However, by applying the same symmetric argument that showed $x \in y$, we can also derive $y \in x$. This means the step where we assumed 
    $y \notin x$ is not justified. With only Axioms 0, 1, and 2, we cannot rule out the case where $x \in y$ and $y \in x$ simultaneously.
\end{proof*}
\end{enumerate}

\end{enumerate}

\end{document}